\section{Resonant modes of a clamped bar}\label{sec:bar}

Resonator loss is a temporal quantity: a standing mode with real wavenumber $k_n=n\pi/L$ decays in time at a complex frequency $\omega_n$. The plane-wave polynomial of \cref{eq:unified} is therefore used with real $k$ and complex $s=-i\omega$.

\subsection{Setting}
Take a bar of length $L$ with both ends clamped ($u=0$) and insulated ($\theta_{,x}=0$). Each eigenfunction has the form $u\propto\sin k_nx$ and $\theta\propto\cos k_nx$ with $k_n=n\pi/L$, so every coupled mode satisfies the boundary conditions exactly and its frequency is a root in $s$ of the polynomial of \cref{eq:unified} at $K=k_n^2$: quartic in $s$ for Lord--Shulman heat flux, cubic for the classical law. The elastic-like root is tracked by continuity in $k$, starting from the adiabatic value $c_1\sqrt{1+\varepsilon}\,k$. In units $c_1=\omega_*=1$, with $\varepsilon=0.1$, the bar-mode figure gives the following.

\subsection{Results}
\begin{itemize}
\item At small $k$ every model gives the same loss, numerically $8.7\times10^{-4}$ at $k=0.01$, about $13\%$ below the first-order Debye value $\varepsilon k/(1+k^2)$ because the temporal and spatial loss measures differ at $O(\varepsilon)$.
\item The classical law gives a Debye-type peak $Q^{-1}=0.0488$ at $k\,c_1/\omega_*=1.07$, and the modal speed falls from the adiabatic value $1.049\,c_1$ to $c_1$ at large $k$.
\item For Lord--Shulman with $\tau\omega_*=0.25$ the peak rises to $0.0676$ at $k=1.39$, and the high-$k$ modal speed settles at $0.984\,c_1$.
\item At $\tau\omega_*=1$, where the thermal front speed equals $c_1$, the peak reaches $0.146$ at $k=1.84$, three times the classical value, and the high-$k$ modal speed is $1.17\,c_1$. This is the degenerate case of equal front speeds. Coupling splits the two fronts as $c^2=c_1^2[1+\varepsilon/2\pm\sqrt{\varepsilon+\varepsilon^2/4}]$, which gives $1.17\,c_1$ for the upper root.
\end{itemize}
The curves are shown in \cref{fig:barmodes}. Near $\tau\omega_*\approx1$ the thermal and elastic branches hybridize, the label ``elastic-like'' loses its sharpness, and results have to be read together with the branch-tracking procedure.

\begin{figure}[t]
\centering
\includegraphics[width=\linewidth]{figs/fig_bar_modes.pdf}
\caption{Loss $Q^{-1}=2|\operatorname{Im}\omega|/|\operatorname{Re}\omega|$ and modal speed of the elastic-like mode of a clamped insulated bar ($\varepsilon=0.1$) against $k=n\pi/L$ for classical and Lord--Shulman heat flux ($\tau\omega_*=0.25$ and $1$).}
\label{fig:barmodes}
\end{figure}

\subsection{Consequence for resonator design}
A bar of fixed length, with its sequence of overtones, traverses the curve of \cref{fig:barmodes} from left to right. The classical model puts the worst loss at $k_nc_1\approx\omega_*$, and the lag raises it only for $\tau\omega_*$ of order one, exactly where the lag law is least trustworthy. For every material of \cref{tab:materials} $\omega_*$ exceeds $10^{11}$\,rad/s. The peak, at $k_nc_1\approx\omega_*$, therefore sits at wavelengths $2\pi c_1/\omega_*$ between about $1$\,nm (fused silica) and $160$\,nm (copper), and micron-sized bar resonators operate on the left, linear-in-$k$ flank where all models agree.

\begin{summarybox}[Summary]
A clamped bar's overtones trace the thermoelastic damping curve from left to right. Classical theory predicts worst loss at $k_nc_1\approx\omega_*$; lagged models raise it only if $\tau\omega_*\sim1$, where the lag law is least trustworthy. For all materials of \cref{tab:materials}, $\omega_*>10^{11}$\,rad/s, so micron-scale bar resonators operate on the linear-in-$k$ flank where all models agree.
\end{summarybox}
